% Einheit 4 von 4 – Ohne Zurücklegen (bank/wahrscheinlichkeit/e4.jsonl, zusammenbau v0.1)
%% TODO Zweigzeile Teil 1: was hier gelernt wird (Satz in Schülersprache aus dem Einheitstitel) – Einheit 4
\einheitenkopf[e4]{Einheit 4 von 4 · Ohne Zurücklegen}
\zweigzeile{ab Kl. 10 · P10 oft}
\setcounter{aufgabe}{20}

%% TODO Titel: Ich-kann-Satz fehlt in der Bank (Platzhalter: Kettenname) – Nr. 21
%% TODO Anweisung: ein Satz über der ersten Teilaufgabe fehlt in der Bank – Nr. 21
\begin{aufgabe}{Ohne Zurücklegen}
%% TODO \rechenplatz{n}: Schrittzahl des Grundfalls fehlt in der Bank (nur bei fester Schreibform, 2.3 b)
\begin{teile}
\teil In einer Dose sind 9 Kekse, 4 davon mit Schoko. Ein Schokokeks wird gegessen. Wie viele Kekse sind noch da, wie viele davon mit Schoko? \leerfeld insgesamt, davon \leerfeld
\teil In einer Urne liegen 3 rote und 5 blaue Kugeln. Zwei Kugeln werden ohne Zurücklegen gezogen. Wie groß ist die Wahrscheinlichkeit für zweimal Rot?
\teil In einem Beutel liegen 4 rote und 3 grüne Kugeln. Zwei werden ohne Zurücklegen gezogen. Ergänze die fehlenden Wahrscheinlichkeiten im Baum.

\baumzwei{rot/\frac{4}{7},grün/}{rot/,grün/\frac{3}{6}}{rot/,grün/}
\teil In einer Urne liegen 3 rote und 2 weiße Kugeln. Zwei werden ohne Zurücklegen gezogen. Wie groß ist die Wahrscheinlichkeit für zwei gleiche Farben?
\teil In einem Hut liegen 7 Zettel für die Aufgaben einer Woche, auf einem steht „frei“. Tom zieht zuerst, dann zieht Lea aus den übrigen Zetteln. Wie groß ist die Wahrscheinlichkeit, dass Lea „frei“ zieht? \hfill (P10 2024 OS)
\teil In einer Box liegen 9 Stifte, 4 davon grün. Drei Stifte werden ohne Zurücklegen gezogen. Wie groß ist die Wahrscheinlichkeit für drei grüne?
\teil In einer Schachtel liegen 15 Pralinen, 5 davon mit Nuss. Emil nimmt nacheinander drei Pralinen, ohne zurückzulegen. Ergänze die beiden leeren Felder im Baum. \hfill (P10 2019 OS)

\baumdrei{Nuss/\frac{5}{15},ohne/}{Nuss/\frac{4}{14},ohne/\frac{10}{14}}{Nuss/\frac{5}{14},ohne/\frac{9}{14}}{Nuss/\frac{3}{13},ohne/\frac{10}{13}}{Nuss/\frac{4}{13},ohne/\frac{9}{13}}{Nuss/,ohne/\frac{9}{13}}{Nuss/\frac{5}{13},ohne/\frac{8}{13}}
\teil Beim Staffellauf wird die Reihenfolge der 9 Läuferinnen zufällig ausgelost; jede läuft genau einmal. Wie groß ist die Wahrscheinlichkeit, dass Pia unter den ersten drei ist? \hfill (P10 2015 OS)
\teil In einer Urne liegen 3 rote und 3 blaue Kugeln. Vergleiche die Wahrscheinlichkeit für zweimal Rot mit und ohne Zurücklegen.
\teil In einer Kiste liegen 15 Stifte: 5 rote, 3 grüne und 7 schwarze. Drei Stifte werden ohne Zurücklegen gezogen. Leo behauptet: „Drei rote sind doppelt so wahrscheinlich wie drei grüne.“ Widerlege die Behauptung mit einer Rechnung. \hfill (P10 2019 OS)
\teil Auf einem Blech liegen 18 Berliner: 15 mit Pflaume, 3 mit Vanille. Ida nimmt ohne hinzusehen nacheinander zwei. Schreibe alle Ergebnisse auf, bei denen mindestens ein Vanille-Berliner dabei ist. Berechne die Wahrscheinlichkeit, dass beide mit Pflaume gefüllt sind. Welches Ereignis gehört zur Rechnung $\frac{3}{18} \cdot \frac{2}{17} + \frac{15}{18} \cdot \frac{14}{17}$? \hfill (P10 2018 OS)
\end{teile}
\end{aufgabe}

%% TODO Titel: Ich-kann-Satz fehlt in der Bank (Platzhalter: Kettenname) – Nr. 22
%% TODO Anweisung: ein Satz über der ersten Teilaufgabe fehlt in der Bank – Nr. 22
\begin{aufgabe}{Ohne Zurücklegen – Fehler finden, Begründen, Darstellungswechsel, Anwendung}
\begin{teile}
\teil In einem Beutel liegen 9 Kugeln, 6 davon blau. Drei werden ohne Zurücklegen gezogen. Nora rechnet für dreimal Blau: $\frac{6}{9} \cdot \frac{5}{8} \cdot \frac{4}{8}$. Finde den Fehler.
\teil Begründe, warum beim zweiten Zug ohne Zurücklegen der Nenner um 1 kleiner ist.
\teil In einem Beutel liegen 5 rote und 3 blaue Kugeln; zwei werden ohne Zurücklegen gezogen. Welches Ereignis gehört zur Rechnung $\frac{5}{8} \cdot \frac{3}{7}$?
\teil In einer Lieferung von 50 Handys sind 2 defekt. Eine Prüferin testet 2 zufällig ausgewählte Handys. Wie groß ist die Wahrscheinlichkeit, dass sie keinen Defekt findet?
\end{teile}
\end{aufgabe}
